% Options for packages loaded elsewhere
% Options for packages loaded elsewhere
\PassOptionsToPackage{unicode}{hyperref}
\PassOptionsToPackage{hyphens}{url}
\PassOptionsToPackage{dvipsnames,svgnames,x11names}{xcolor}
%
\documentclass[
  10pt,
  letterpaper,
  DIV=11,
  numbers=noendperiod]{scrartcl}
\usepackage{xcolor}
\usepackage{amsmath,amssymb}
\setcounter{secnumdepth}{-\maxdimen} % remove section numbering
\usepackage{iftex}
\ifPDFTeX
  \usepackage[T1]{fontenc}
  \usepackage[utf8]{inputenc}
  \usepackage{textcomp} % provide euro and other symbols
\else % if luatex or xetex
  \usepackage{unicode-math} % this also loads fontspec
  \defaultfontfeatures{Scale=MatchLowercase}
  \defaultfontfeatures[\rmfamily]{Ligatures=TeX,Scale=1}
\fi
\usepackage{lmodern}
\ifPDFTeX\else
  % xetex/luatex font selection
  \setmonofont[]{Fira Code}
  \setmathfont[]{TeX Gyre Pagella Math}
\fi
% Use upquote if available, for straight quotes in verbatim environments
\IfFileExists{upquote.sty}{\usepackage{upquote}}{}
\IfFileExists{microtype.sty}{% use microtype if available
  \usepackage[]{microtype}
  \UseMicrotypeSet[protrusion]{basicmath} % disable protrusion for tt fonts
}{}
\usepackage{setspace}
\makeatletter
\@ifundefined{KOMAClassName}{% if non-KOMA class
  \IfFileExists{parskip.sty}{%
    \usepackage{parskip}
  }{% else
    \setlength{\parindent}{0pt}
    \setlength{\parskip}{6pt plus 2pt minus 1pt}}
}{% if KOMA class
  \KOMAoptions{parskip=half}}
\makeatother
% Make \paragraph and \subparagraph free-standing
\makeatletter
\ifx\paragraph\undefined\else
  \let\oldparagraph\paragraph
  \renewcommand{\paragraph}{
    \@ifstar
      \xxxParagraphStar
      \xxxParagraphNoStar
  }
  \newcommand{\xxxParagraphStar}[1]{\oldparagraph*{#1}\mbox{}}
  \newcommand{\xxxParagraphNoStar}[1]{\oldparagraph{#1}\mbox{}}
\fi
\ifx\subparagraph\undefined\else
  \let\oldsubparagraph\subparagraph
  \renewcommand{\subparagraph}{
    \@ifstar
      \xxxSubParagraphStar
      \xxxSubParagraphNoStar
  }
  \newcommand{\xxxSubParagraphStar}[1]{\oldsubparagraph*{#1}\mbox{}}
  \newcommand{\xxxSubParagraphNoStar}[1]{\oldsubparagraph{#1}\mbox{}}
\fi
\makeatother

\usepackage{color}
\usepackage{fancyvrb}
\newcommand{\VerbBar}{|}
\newcommand{\VERB}{\Verb[commandchars=\\\{\}]}
\DefineVerbatimEnvironment{Highlighting}{Verbatim}{commandchars=\\\{\}}
% Add ',fontsize=\small' for more characters per line
\usepackage{framed}
\definecolor{shadecolor}{RGB}{241,243,245}
\newenvironment{Shaded}{\begin{snugshade}}{\end{snugshade}}
\newcommand{\AlertTok}[1]{\textcolor[rgb]{0.68,0.00,0.00}{#1}}
\newcommand{\AnnotationTok}[1]{\textcolor[rgb]{0.37,0.37,0.37}{#1}}
\newcommand{\AttributeTok}[1]{\textcolor[rgb]{0.40,0.45,0.13}{#1}}
\newcommand{\BaseNTok}[1]{\textcolor[rgb]{0.68,0.00,0.00}{#1}}
\newcommand{\BuiltInTok}[1]{\textcolor[rgb]{0.00,0.23,0.31}{#1}}
\newcommand{\CharTok}[1]{\textcolor[rgb]{0.13,0.47,0.30}{#1}}
\newcommand{\CommentTok}[1]{\textcolor[rgb]{0.37,0.37,0.37}{#1}}
\newcommand{\CommentVarTok}[1]{\textcolor[rgb]{0.37,0.37,0.37}{\textit{#1}}}
\newcommand{\ConstantTok}[1]{\textcolor[rgb]{0.56,0.35,0.01}{#1}}
\newcommand{\ControlFlowTok}[1]{\textcolor[rgb]{0.00,0.23,0.31}{\textbf{#1}}}
\newcommand{\DataTypeTok}[1]{\textcolor[rgb]{0.68,0.00,0.00}{#1}}
\newcommand{\DecValTok}[1]{\textcolor[rgb]{0.68,0.00,0.00}{#1}}
\newcommand{\DocumentationTok}[1]{\textcolor[rgb]{0.37,0.37,0.37}{\textit{#1}}}
\newcommand{\ErrorTok}[1]{\textcolor[rgb]{0.68,0.00,0.00}{#1}}
\newcommand{\ExtensionTok}[1]{\textcolor[rgb]{0.00,0.23,0.31}{#1}}
\newcommand{\FloatTok}[1]{\textcolor[rgb]{0.68,0.00,0.00}{#1}}
\newcommand{\FunctionTok}[1]{\textcolor[rgb]{0.28,0.35,0.67}{#1}}
\newcommand{\ImportTok}[1]{\textcolor[rgb]{0.00,0.46,0.62}{#1}}
\newcommand{\InformationTok}[1]{\textcolor[rgb]{0.37,0.37,0.37}{#1}}
\newcommand{\KeywordTok}[1]{\textcolor[rgb]{0.00,0.23,0.31}{\textbf{#1}}}
\newcommand{\NormalTok}[1]{\textcolor[rgb]{0.00,0.23,0.31}{#1}}
\newcommand{\OperatorTok}[1]{\textcolor[rgb]{0.37,0.37,0.37}{#1}}
\newcommand{\OtherTok}[1]{\textcolor[rgb]{0.00,0.23,0.31}{#1}}
\newcommand{\PreprocessorTok}[1]{\textcolor[rgb]{0.68,0.00,0.00}{#1}}
\newcommand{\RegionMarkerTok}[1]{\textcolor[rgb]{0.00,0.23,0.31}{#1}}
\newcommand{\SpecialCharTok}[1]{\textcolor[rgb]{0.37,0.37,0.37}{#1}}
\newcommand{\SpecialStringTok}[1]{\textcolor[rgb]{0.13,0.47,0.30}{#1}}
\newcommand{\StringTok}[1]{\textcolor[rgb]{0.13,0.47,0.30}{#1}}
\newcommand{\VariableTok}[1]{\textcolor[rgb]{0.07,0.07,0.07}{#1}}
\newcommand{\VerbatimStringTok}[1]{\textcolor[rgb]{0.13,0.47,0.30}{#1}}
\newcommand{\WarningTok}[1]{\textcolor[rgb]{0.37,0.37,0.37}{\textit{#1}}}

\usepackage{longtable,booktabs,array}
\usepackage{calc} % for calculating minipage widths
% Correct order of tables after \paragraph or \subparagraph
\usepackage{etoolbox}
\makeatletter
\patchcmd\longtable{\par}{\if@noskipsec\mbox{}\fi\par}{}{}
\makeatother
% Allow footnotes in longtable head/foot
\IfFileExists{footnotehyper.sty}{\usepackage{footnotehyper}}{\usepackage{footnote}}
\makesavenoteenv{longtable}
\usepackage{graphicx}
\makeatletter
\newsavebox\pandoc@box
\newcommand*\pandocbounded[1]{% scales image to fit in text height/width
  \sbox\pandoc@box{#1}%
  \Gscale@div\@tempa{\textheight}{\dimexpr\ht\pandoc@box+\dp\pandoc@box\relax}%
  \Gscale@div\@tempb{\linewidth}{\wd\pandoc@box}%
  \ifdim\@tempb\p@<\@tempa\p@\let\@tempa\@tempb\fi% select the smaller of both
  \ifdim\@tempa\p@<\p@\scalebox{\@tempa}{\usebox\pandoc@box}%
  \else\usebox{\pandoc@box}%
  \fi%
}
% Set default figure placement to htbp
\def\fps@figure{htbp}
\makeatother




\pagestyle{plain}

\setlength{\emergencystretch}{3em} % prevent overfull lines

\providecommand{\tightlist}{%
  \setlength{\itemsep}{0pt}\setlength{\parskip}{0pt}}



 


%\usepackage{amsfonts} % 
\usepackage{amssymb}
\usepackage{mathtools} % For \DeclarePairedDelimiter
\usepackage{multicol} % For multiple columns
\usepackage{graphicx} % For images
\usepackage{float} % For figure placement
% ==========================
% Sets & Blackboard bold
% ==========================
\newcommand{\RR}{\mathbb{R}}
\newcommand{\NN}{\mathbb{N}}
\newcommand{\PP}{\mathbb{P}}
\newcommand{\QQ}{\mathbb{Q}}
\newcommand{\ZZ}{\mathbb{Z}}
% ==========================
% Math macros
% ==========================
% Expectation & Probability
% ==========================
\newcommand{\E}[1]{\mathbb{E}\left[ #1 \right]}
\newcommand{\Es}[2]{\mathbb{E}_{#1}\left[ #2 \right]}  % subscripted expectation
\newcommand{\Prob}[1]{\mathbb{P}\left( #1 \right)}     % probability with parentheses
\newcommand{\pr}[1]{\PP\!\left( #1 \right)}
\newcommand{\EE}[1]{\E\!\left[ #1 \right]}            % alias for expectation
\newcommand{\Ehat}[1]{\hat{\mathbb{E}}\left[#1\right]}
% ==========================
% Propensity symbols (optional)
% ==========================
\newcommand{\PPhi}{\mathcal{P}}       % propensity score symbol
\newcommand{\PPhiHat}{\hat{\mathcal{P}}} % estimated propensity

% ==========================
% Vectors
% ==========================
\newcommand{\vect}[1]{\mathbf{#1}}

% ==========================
% Operators (variance/covariance/correlation/logit/sgn)
% ==========================
% Canonical operator names (upright, proper spacing)
\DeclareMathOperator{\Var}{Var}
\DeclareMathOperator{\Cov}{Cov}
\DeclareMathOperator{\Cor}{Cor}
\DeclareMathOperator{\logit}{logit}
\DeclareMathOperator{\sgn}{sgn}
% Convenience wrappers
\newcommand{\var}[1]{\Var\!\left( #1 \right)}
\newcommand{\cov}[2]{\Cov\!\left( #1,\, #2 \right)}
\newcommand{\cor}[2]{\Cor\!\left( #1,\, #2 \right)}

% ==========================
% Norms, absolute value, inner product (paired delimiters)
% ==========================
\DeclarePairedDelimiter{\abs}{\lvert}{\rvert}
\DeclarePairedDelimiter{\norm}{\lVert}{\rVert}
\DeclarePairedDelimiter{\ip}{\langle}{\rangle} % inner product
% Usage: \abs{x}, \norm{x}. For auto-sized delimiters: \abs*{x}, \norm*{x}.

% ==========================
% Indicator (two variants)
% ==========================
\newcommand{\indf}[1]{\mathbf{1}\left\{ #1 \right\}} % bold 1 style
\newcommand{\ind}[1]{\mathbb{I}\left\{ #1 \right\}}  % double-struck I style
\newcommand{\I}[1]{\ind{#1}}                            % alias

% ==========================
% Independence / Non-independence
%  - Bare symbol:  $X \indep Y$, $X \nindep Y$, $X \notindep Y$
%  - Braced args:  $\indep{X}{Y}$, $\nindep{X}{Y}$, $\notindep{X}{Y}$
% ==========================
\makeatletter
\newcommand{\indep}{\@ifnextchar\bgroup\indep@two\indep@bare}
\newcommand{\indep@two}[2]{#1 \, \mathrel{\perp\!\!\!\perp} \, #2}
\newcommand{\indep@bare}{\mathrel{\perp\!\!\!\perp}}

\newcommand{\nindep}{\@ifnextchar\bgroup\nindep@two\nindep@bare}
\newcommand{\nindep@two}[2]{#1 \, \not\!\perp\!\!\!\perp \, #2}
\newcommand{\nindep@bare}{\not\!\perp\!\!\!\perp}

\newcommand{\notindep}{\@ifnextchar\bgroup\notindep@two\notindep@bare}
\newcommand{\notindep@two}[2]{#1 \, \not\!\perp\!\!\!\perp \, #2}
\newcommand{\notindep@bare}{\not\!\perp\!\!\!\perp}
\makeatother
% ==========================
% Interventions / Graphical
% ==========================
\newcommand{\doop}[1]{\mathrm{do}\!\left( #1 \right)}  % do-operator
\newcommand{\dsep}{\mathrel{\perp\!\!\!\perp}}       % d-separation relation

% ==========================
% Potential outcomes & shorthands
% ==========================
\newcommand{\Y}[1]{Y^{#1}}                 % Y^a
\newcommand{\Ya}{Y^a}
\newcommand{\Yone}{Y^1}
\newcommand{\Yzero}{Y^0}

% ==========================
% Standard causal estimands
% ==========================
\newcommand{\DefATE}{\E\!\left[ \Yone - \Yzero \right]}
\newcommand{\DefATT}{\E\!\left[ \Yone - \Yzero \mid A=1 \right]}
\newcommand{\DefATC}{\E\!\left[ \Yone - \Yzero \mid A=0 \right]}
\newcommand{\DefCATE}[1]{\E\!\left[ \Yone - \Yzero \mid #1 \right]}

% ==========================
% Identification building blocks
% ==========================
\newcommand{\gform}[1]{\sum_{l} \E\!\left[ Y \mid A=#1,\, L=l \right] \, \PP(L=l)}
\newcommand{\ipw}[2]{\frac{\indf{#1}}{\PP(#1\mid #2)}} % stabilized IPW form

% ==========================
% Risk measures
% ==========================
\DeclareMathOperator{\RD}{RD}
\DeclareMathOperator{\OR}{OR}

% ==========================
% Optimization operators
% ==========================
\DeclareMathOperator*{\argmin}{arg\,min}
\DeclareMathOperator*{\argmax}{arg\,max}
\KOMAoption{captions}{tableheading}
\makeatletter
\@ifpackageloaded{caption}{}{\usepackage{caption}}
\AtBeginDocument{%
\ifdefined\contentsname
  \renewcommand*\contentsname{Table of contents}
\else
  \newcommand\contentsname{Table of contents}
\fi
\ifdefined\listfigurename
  \renewcommand*\listfigurename{List of Figures}
\else
  \newcommand\listfigurename{List of Figures}
\fi
\ifdefined\listtablename
  \renewcommand*\listtablename{List of Tables}
\else
  \newcommand\listtablename{List of Tables}
\fi
\ifdefined\figurename
  \renewcommand*\figurename{Figure}
\else
  \newcommand\figurename{Figure}
\fi
\ifdefined\tablename
  \renewcommand*\tablename{Table}
\else
  \newcommand\tablename{Table}
\fi
}
\@ifpackageloaded{float}{}{\usepackage{float}}
\floatstyle{ruled}
\@ifundefined{c@chapter}{\newfloat{codelisting}{h}{lop}}{\newfloat{codelisting}{h}{lop}[chapter]}
\floatname{codelisting}{Listing}
\newcommand*\listoflistings{\listof{codelisting}{List of Listings}}
\makeatother
\makeatletter
\makeatother
\makeatletter
\@ifpackageloaded{caption}{}{\usepackage{caption}}
\@ifpackageloaded{subcaption}{}{\usepackage{subcaption}}
\makeatother
\usepackage{bookmark}
\IfFileExists{xurl.sty}{\usepackage{xurl}}{} % add URL line breaks if available
\urlstyle{same}
\hypersetup{
  pdftitle={Fundamental Concepts in Causal Inference},
  pdfauthor={Davood Tofighi, Ph.D.},
  colorlinks=true,
  linkcolor={blue},
  filecolor={Maroon},
  citecolor={Blue},
  urlcolor={Blue},
  pdfcreator={LaTeX via pandoc}}


\title{Fundamental Concepts in Causal Inference}
\author{Davood Tofighi, Ph.D.}
\date{}
\begin{document}
\maketitle


\setstretch{1.25}
Welcome. Let's begin our exploration of the fundamental concepts that
form the bedrock of causal inference. Before we can navigate the
complexities of methods like IP weighting or g-estimation, we must have
a rock-solid understanding of the basic epidemiological measures used to
quantify associations and, under specific assumptions, causal effects.

Think of these concepts as the essential vocabulary of our field.
Mastering them will allow you to precisely state your causal question,
understand the measures reported in the scientific literature, and
correctly interpret the output of your own analyses. We will draw
heavily from Chapters 1 and 17 of Hernán and Robins (HR) to build this
foundation.

Let's start with the most critical distinction of all.

\begin{center}\rule{0.5\linewidth}{0.5pt}\end{center}

\subsection{The Core Distinction: Association
vs.~Causation}\label{the-core-distinction-association-vs.-causation}

👨‍🏫 \textbf{Intuitive Understanding:}

At its heart, epidemiology is often concerned with figuring out if an
\textbf{exposure} (like a medication, a lifestyle choice, or an
environmental factor) causes an \textbf{outcome} (like a disease,
recovery, or death).

An \textbf{association} is simply an observed statistical relationship.
We observe that people with the exposure are more or less likely to have
the outcome than people without the exposure. For example, we might
observe that individuals who carry a lighter in their pocket are more
likely to develop lung cancer than those who do not. This is a real,
observable association.

A \textbf{causal effect}, however, is a much stronger claim. It means
that the exposure \emph{itself} influences the outcome. For the lighter
example, we intuitively know the lighter doesn't cause cancer. The
association exists because a third factor, \textbf{smoking}, is a common
cause of both carrying a lighter and developing lung cancer.

This simple example illustrates the classic adage: \textbf{``Association
is not causation.''}

Our central task in this course is to determine the conditions under
which an observed association can be interpreted as a causal effect. To
do this, we need a formal language.

\begin{center}\rule{0.5\linewidth}{0.5pt}\end{center}

\subsubsection{Mathematical Foundations}\label{mathematical-foundations}

The language we use is that of \textbf{potential outcomes} (or
counterfactuals), as introduced in HR Chapter 1{[}cite: 38{]}. For an
individual \(i\), and a binary exposure \(A\) (1=exposed, 0=unexposed),
we define two potential outcomes:

\begin{itemize}
\tightlist
\item
  \(Y_i^{a=1}\): The outcome for individual \(i\) \emph{if} they were
  exposed.
\item
  \(Y_i^{a=0}\): The outcome for individual \(i\) \emph{if} they were
  unexposed.
\end{itemize}

\textbf{Causation} is about comparing these potential outcomes \emph{for
the same individuals}. A causal effect exists for an individual if
\(Y_i^{a=1} \neq Y_i^{a=0}\){[}cite: 38{]}. Since we can't observe both
potential outcomes for the same person at the same time (this is the
\textbf{fundamental problem of causal inference}), we focus on the
\textbf{average causal effect (ATE)} in a population:

\[
\text{Average Causal Effect} \iff E[Y^{a=1}] \neq E[Y^{a=0}]
\]

\textbf{Association}, on the other hand, is about comparing observed
outcomes in different groups of people{[}cite: 96{]}.

\[
\text{Association} \iff E[Y|A=1] \neq E[Y|A=0]
\]

The crucial link between these two concepts is the assumption of
\textbf{exchangeability}, which we will discuss extensively. For now,
understand that if the exposed and unexposed groups are exchangeable
(i.e., they are comparable in every way except for the exposure), then
the association we measure will equal the causal effect we seek. An
ideal randomized trial achieves this. Our job with observational data is
to try to achieve it through careful study design and analysis.

\begin{center}\rule{0.5\linewidth}{0.5pt}\end{center}

\subsection{Quantifying Association and Causal
Effects}\label{quantifying-association-and-causal-effects}

Whether we are describing an association or a causal effect, we use the
same set of epidemiological measures. The only difference is the
notation and, most importantly, the assumptions that justify the
interpretation.

For the following, let's consider a binary outcome \(Y\) (1=event, 0=no
event), where the mean \(E[Y]\) is simply the probability of the event,
\(Pr(Y=1)\), which we often call the \textbf{risk}.

\subsubsection{Risk Difference (RD)}\label{risk-difference-rd}

\begin{itemize}
\tightlist
\item
  \textbf{Intuitive Understanding:} The RD tells us the \textbf{absolute
  difference} in risk between two groups. It answers the question: ``How
  many more (or fewer) cases of the outcome are there in the exposed
  group compared to the unexposed group, out of every 100 people?'' It's
  a useful measure for public health because it directly translates to
  the number of cases that could be prevented.
\item
  \textbf{Mathematical Foundations:}

  \begin{itemize}
  \tightlist
  \item
    \textbf{Associational Risk Difference:}
    \(RD_{Assoc} = Pr(Y=1|A=1) - Pr(Y=1|A=0)\) {[}cite: 84{]}
  \item
    \textbf{Causal Risk Difference:}
    \(RD_{Causal} = Pr(Y^{a=1}=1) - Pr(Y^{a=0}=1)\) {[}cite: 57{]}
  \end{itemize}
\end{itemize}

\subsubsection{Risk Ratio (RR)}\label{risk-ratio-rr}

\begin{itemize}
\tightlist
\item
  \textbf{Intuitive Understanding:} The RR, also called relative risk,
  tells us the \textbf{relative difference} in risk. It answers the
  question: ``How many times more (or less) likely are exposed
  individuals to develop the outcome compared to unexposed
  individuals?'' It's often used to measure the strength of an
  association. An RR of 2 means the risk is ``doubled,'' while an RR of
  0.5 means the risk is ``halved.''
\item
  \textbf{Mathematical Foundations:}

  \begin{itemize}
  \tightlist
  \item
    \textbf{Associational Risk Ratio:}
    \(RR_{Assoc} = \frac{Pr(Y=1|A=1)}{Pr(Y=1|A=0)}\) {[}cite: 84{]}
  \item
    \textbf{Causal Risk Ratio:}
    \(RR_{Causal} = \frac{Pr(Y^{a=1}=1)}{Pr(Y^{a=0}=1)}\) {[}cite: 57{]}
  \end{itemize}
\end{itemize}

\subsubsection{Odds And the Odds Ratio
(OR)}\label{odds-and-the-odds-ratio-or}

\begin{itemize}
\tightlist
\item
  \textbf{Intuitive Understanding:} This one is less intuitive. The
  \textbf{odds} of an event is the probability of it happening divided
  by the probability of it not happening. If the probability of an event
  is 0.75 (or 75\%), the odds are \(0.75 / (1-0.75) = 3\), often stated
  as ``3 to 1.'' The \textbf{Odds Ratio (OR)} is simply the ratio of the
  odds in the exposed group to the odds in the unexposed group. A key
  property is that when the outcome is rare, the OR is a good
  approximation of the RR.
\item
  \textbf{Mathematical Foundations:}

  \begin{itemize}
  \tightlist
  \item
    The odds of an event is
    \(Odds = \frac{Pr(Y=1)}{Pr(Y=0)} = \frac{p}{1-p}\).
  \item
    \textbf{Associational Odds Ratio:}
    \(OR_{Assoc} = \frac{Pr(Y=1|A=1) / Pr(Y=0|A=1)}{Pr(Y=1|A=0) / Pr(Y=0|A=0)}\)
    {[}cite: 85{]}
  \item
    \textbf{Causal Odds Ratio:}
    \(OR_{Causal} = \frac{Pr(Y^{a=1}=1) / Pr(Y^{a=1}=0)}{Pr(Y^{a=0}=1) / Pr(Y^{a=0}=0)}\)
    {[}cite: 57{]}
  \end{itemize}
\end{itemize}

\begin{center}\rule{0.5\linewidth}{0.5pt}\end{center}

\subsection{Measures for Survival
Analysis}\label{measures-for-survival-analysis}

When our outcome is the \textbf{time until an event occurs} (e.g.,
death, disease diagnosis), we use slightly different measures, as
detailed in HR Chapter 17{[}cite: 892{]}.

\subsubsection{\texorpdfstring{\textbf{Survival Probability and Risk
(Cumulative
Incidence)}}{Survival Probability and Risk (Cumulative Incidence)}}\label{survival-probability-and-risk-cumulative-incidence}

\begin{itemize}
\tightlist
\item
  \textbf{Intuitive Understanding:} The \textbf{survival probability} at
  time \(k\), denoted \(S(k)\), is the probability of not having the
  event by that time (i.e., ``surviving'' past time \(k\)). A survival
  curve plots this probability over time. The \textbf{risk}, or
  \textbf{cumulative incidence}, at time \(k\), denoted \(F(k)\), is the
  opposite: the probability of having had the event by that time.
\item
  \textbf{Mathematical Foundations:}

  \begin{itemize}
  \tightlist
  \item
    Let \(T\) be the time to the event.
  \item
    \textbf{Survival Probability:} \(S(k) = Pr(T \> k)\) {[}cite: 893{]}
  \item
    \textbf{Risk / Cumulative Incidence:}
    \(F(k) = Pr(T \le k) = 1 - S(k)\) {[}cite: 893{]}
  \end{itemize}
\end{itemize}

\subsubsection{\texorpdfstring{\textbf{Hazard}}{Hazard}}\label{hazard}

\begin{itemize}
\tightlist
\item
  \textbf{Intuitive Understanding:} The \textbf{hazard} is the
  instantaneous potential for the event to occur at time \(k\), given
  that an individual has survived up to that point. Think of it as the
  risk of the event happening \emph{right now}, among those who are
  still at risk. Unlike cumulative risk which can only increase, the
  hazard can go up or down over time.
\item
  \textbf{Mathematical Foundations:}

  \begin{itemize}
  \tightlist
  \item
    For discrete time (e.g., months), the hazard at month \(k\) is the
    conditional probability: \(h(k) = Pr(T=k | T \> k-1)\){[}cite:
    895{]}.
  \end{itemize}
\end{itemize}

These measures form the basis for comparing groups in survival analysis,
often through survival curves or hazard ratios.

This concludes our foundational overview. The next crucial step is to
understand the assumptions that allow us to move from measuring
associations in our data to estimating these causal quantities.

\begin{center}\rule{0.5\linewidth}{0.5pt}\end{center}

\subsection{Appendix: Comprehension and
Application}\label{appendix-comprehension-and-application}

\subsubsection{Part A: Conceptual
Questions}\label{part-a-conceptual-questions}

\begin{enumerate}
\def\labelenumi{\arabic{enumi}.}
\tightlist
\item
  \textbf{The Adage:} Explain in your own words, using a simple public
  health example (other than lighters/smoking), why ``association is not
  causation.'' Draw a simple causal diagram for your example.
\item
  \textbf{Choosing a Measure:} An intervention reduces the risk of a
  disease from 4 in 100 to 2 in 100. Another intervention reduces the
  risk from 40 in 100 to 30 in 100.

  \begin{itemize}
  \tightlist
  \item
    Calculate the Risk Ratio (RR) and Risk Difference (RD) for both
    interventions.
  \item
    Which measure (RR or RD) better reflects the public health impact
    (number of cases prevented)? Which measure better reflects the
    strength of the effect relative to the baseline risk?
  \end{itemize}
\item
  \textbf{Counterfactuals:} A patient took an aspirin and their headache
  went away. Why is it impossible to know for sure if the aspirin had a
  causal effect on that specific patient's headache? What is the
  specific term for this issue?
\end{enumerate}

\subsubsection{Part B: Computation from a
Table}\label{part-b-computation-from-a-table}

Consider the following data from a hypothetical study of 2000
individuals on the effect of a new blood pressure drug (\(A\)) on the
one-year risk of stroke (\(Y\)).

\begin{longtable}[]{@{}llll@{}}
\toprule\noalign{}
& Stroke (\(Y=1\)) & No Stroke (\(Y=0\)) & Total \\
\midrule\noalign{}
\endhead
\bottomrule\noalign{}
\endlastfoot
\textbf{Drug (\(A=1\))} & 80 & 920 & 1000 \\
\textbf{No Drug (\(A=0\))} & 50 & 950 & 1000 \\
\end{longtable}

Calculate and interpret the following \textbf{associational} measures:

\begin{enumerate}
\def\labelenumi{\arabic{enumi}.}
\tightlist
\item
  Risk of stroke in the treated and untreated groups.
\item
  Risk Difference (RD).
\item
  Risk Ratio (RR).
\item
  Odds Ratio (OR).
\end{enumerate}

\subsubsection{Part C: Application with
R}\label{part-c-application-with-r}

Let's simulate a small dataset to see how this works in R.

\textbf{Scenario:} We will simulate a population where a treatment
\texttt{A} has a true causal effect on a binary outcome \texttt{Y}.
We'll start with a simple case: a randomized trial where there is no
confounding.

\begin{Shaded}
\begin{Highlighting}[]
\CommentTok{\# Load the tidyverse package for data manipulation and plotting}
\CommentTok{\# If you don\textquotesingle{}t have it, run: install.packages("tidyverse")}
\FunctionTok{library}\NormalTok{(tidyverse)}

\CommentTok{\# {-}{-}{-} 1. Simulate a simple Randomized Controlled Trial (RCT) {-}{-}{-}}
\FunctionTok{set.seed}\NormalTok{(}\DecValTok{579}\NormalTok{) }\CommentTok{\# for reproducibility}

\NormalTok{n }\OtherTok{\textless{}{-}} \DecValTok{2000} \CommentTok{\# number of individuals in our study}
\CommentTok{\# In an RCT, treatment A is assigned randomly (e.g., with 50\% probability)}
\CommentTok{\# It does not depend on any other characteristic of the individual.}
\NormalTok{treatment }\OtherTok{\textless{}{-}} \FunctionTok{rbinom}\NormalTok{(n, }\DecValTok{1}\NormalTok{, }\FloatTok{0.5}\NormalTok{)}

\CommentTok{\# The baseline risk of the outcome Y is 10\%}
\CommentTok{\# The true causal effect of the treatment is to increase the probability by 15 percentage points}
\CommentTok{\# Pr(Y=1 | A=0) = 0.10}
\CommentTok{\# Pr(Y=1 | A=1) = 0.10 + 0.15 = 0.25}
\CommentTok{\# The true Causal Risk Ratio is 0.25 / 0.10 = 2.5}
\CommentTok{\# The true Causal Risk Difference is 0.25 {-} 0.10 = 0.15}
\NormalTok{prob\_y }\OtherTok{\textless{}{-}} \FunctionTok{ifelse}\NormalTok{(treatment }\SpecialCharTok{==} \DecValTok{1}\NormalTok{, }\FloatTok{0.25}\NormalTok{, }\FloatTok{0.10}\NormalTok{)}
\NormalTok{outcome }\OtherTok{\textless{}{-}} \FunctionTok{rbinom}\NormalTok{(n, }\DecValTok{1}\NormalTok{, prob\_y)}

\CommentTok{\# Create a data frame (our dataset)}
\NormalTok{rct\_data }\OtherTok{\textless{}{-}} \FunctionTok{data.frame}\NormalTok{(}\AttributeTok{A =}\NormalTok{ treatment, }\AttributeTok{Y =}\NormalTok{ outcome)}

\CommentTok{\# Let\textquotesingle{}s inspect the first few rows}
\FunctionTok{head}\NormalTok{(rct\_data)}

\CommentTok{\# {-}{-}{-} 2. Calculate Associational Measures from the Data {-}{-}{-}}

\CommentTok{\# We can create a 2x2 table to summarize the data}
\NormalTok{rct\_table }\OtherTok{\textless{}{-}} \FunctionTok{table}\NormalTok{(rct\_data}\SpecialCharTok{$}\NormalTok{A, rct\_data}\SpecialCharTok{$}\NormalTok{Y)}
\FunctionTok{print}\NormalTok{(}\StringTok{"2x2 Table (Rows: Treatment A, Columns: Outcome Y)"}\NormalTok{)}
\FunctionTok{print}\NormalTok{(rct\_table)}

\CommentTok{\# Let\textquotesingle{}s calculate the measures}
\NormalTok{summary\_stats }\OtherTok{\textless{}{-}}\NormalTok{ rct\_data }\SpecialCharTok{\%\textgreater{}\%}
  \FunctionTok{group\_by}\NormalTok{(A) }\SpecialCharTok{\%\textgreater{}\%}
  \FunctionTok{summarise}\NormalTok{(}
    \AttributeTok{n =} \FunctionTok{n}\NormalTok{(),}
    \AttributeTok{events =} \FunctionTok{sum}\NormalTok{(Y),}
    \AttributeTok{risk =} \FunctionTok{mean}\NormalTok{(Y) }\CommentTok{\# For a binary 0/1 outcome, mean() is the same as risk}
\NormalTok{  ) }\SpecialCharTok{\%\textgreater{}\%}
  \FunctionTok{ungroup}\NormalTok{()}

\FunctionTok{print}\NormalTok{(}\StringTok{"Summary Statistics by Treatment Group:"}\NormalTok{)}
\FunctionTok{print}\NormalTok{(summary\_stats)}

\CommentTok{\# Extract risks for each group}
\NormalTok{risk\_treated }\OtherTok{\textless{}{-}}\NormalTok{ summary\_stats}\SpecialCharTok{$}\NormalTok{risk[summary\_stats}\SpecialCharTok{$}\NormalTok{A }\SpecialCharTok{==} \DecValTok{1}\NormalTok{]}
\NormalTok{risk\_untreated }\OtherTok{\textless{}{-}}\NormalTok{ summary\_stats}\SpecialCharTok{$}\NormalTok{risk[summary\_stats}\SpecialCharTok{$}\NormalTok{A }\SpecialCharTok{==} \DecValTok{0}\NormalTok{]}

\CommentTok{\# Calculate Associational RD and RR}
\NormalTok{assoc\_rd }\OtherTok{\textless{}{-}}\NormalTok{ risk\_treated }\SpecialCharTok{{-}}\NormalTok{ risk\_untreated}
\NormalTok{assoc\_rr }\OtherTok{\textless{}{-}}\NormalTok{ risk\_treated }\SpecialCharTok{/}\NormalTok{ risk\_untreated}

\CommentTok{\# Calculate Odds Ratio}
\NormalTok{odds\_treated }\OtherTok{\textless{}{-}}\NormalTok{ risk\_treated }\SpecialCharTok{/}\NormalTok{ (}\DecValTok{1} \SpecialCharTok{{-}}\NormalTok{ risk\_treated)}
\NormalTok{odds\_untreated }\OtherTok{\textless{}{-}}\NormalTok{ risk\_untreated }\SpecialCharTok{/}\NormalTok{ (}\DecValTok{1} \SpecialCharTok{{-}}\NormalTok{ risk\_untreated)}
\NormalTok{assoc\_or }\OtherTok{\textless{}{-}}\NormalTok{ odds\_treated }\SpecialCharTok{/}\NormalTok{ odds\_untreated}

\CommentTok{\# {-}{-}{-} 3. Interpretation {-}{-}{-}}

\FunctionTok{cat}\NormalTok{(}\StringTok{"}\SpecialCharTok{\textbackslash{}n}\StringTok{{-}{-}{-} Results from our Simulated RCT {-}{-}{-}}\SpecialCharTok{\textbackslash{}n}\StringTok{"}\NormalTok{)}
\FunctionTok{cat}\NormalTok{(}\StringTok{"Associational Risk Difference:"}\NormalTok{, }\FunctionTok{round}\NormalTok{(assoc\_rd, }\DecValTok{3}\NormalTok{), }\StringTok{"}\SpecialCharTok{\textbackslash{}n}\StringTok{"}\NormalTok{)}
\FunctionTok{cat}\NormalTok{(}\StringTok{"Associational Risk Ratio:"}\NormalTok{, }\FunctionTok{round}\NormalTok{(assoc\_rr, }\DecValTok{3}\NormalTok{), }\StringTok{"}\SpecialCharTok{\textbackslash{}n}\StringTok{"}\NormalTok{)}
\FunctionTok{cat}\NormalTok{(}\StringTok{"Associational Odds Ratio:"}\NormalTok{, }\FunctionTok{round}\NormalTok{(assoc\_or, }\DecValTok{3}\NormalTok{), }\StringTok{"}\SpecialCharTok{\textbackslash{}n}\StringTok{"}\NormalTok{)}
\end{Highlighting}
\end{Shaded}

\textbf{R Syntax Explanation:}

\begin{itemize}
\tightlist
\item
  \texttt{set.seed(579)}: This function ensures that every time you run
  the code, the ``random'' numbers generated will be the same. This is
  crucial for making our results reproducible.
\item
  \texttt{rbinom(n,\ 1,\ 0.5)}: This generates \texttt{n} random numbers
  from a binomial distribution. Here, it's like flipping a coin
  \texttt{n} times, where the probability of getting a 1 (heads) is 0.5.
  This simulates our random treatment assignment.
\item
  \texttt{ifelse(test,\ yes,\ no)}: This is a conditional statement. It
  checks if \texttt{treatment\ ==\ 1}. If true, it assigns the value
  \texttt{0.25}; if false, it assigns \texttt{0.10}. This is how we
  build the true causal effect into our simulation.
\item
  \texttt{\%\textgreater{}\%}: This is the ``pipe'' operator from the
  \texttt{dplyr} package (part of \texttt{tidyverse}). It takes the
  output of the expression on its left and ``pipes'' it in as the first
  argument to the function on its right. It makes the code easier to
  read.
\item
  \texttt{group\_by(A)}: This tells R to perform the next set of
  operations separately for each group defined by the variable
  \texttt{A} (treated and untreated).
\item
  \texttt{summarise(…)}: This creates a new data frame with summary
  statistics. \texttt{n()} counts the rows, \texttt{sum(Y)} counts the
  events, and \texttt{mean(Y)} calculates the risk.
\item
  \texttt{cat(…)}: This function is like \texttt{print()}, but gives us
  more control over formatting the output text. The
  \texttt{\textbackslash{}n} creates a new line.
\end{itemize}

\textbf{Interpretation of R Output:}

When you run this code, you will see that the associational measures you
calculated from the simulated data are very close to the \emph{true
causal effects} we built into the simulation (Associational RD ≈ 0.15,
Associational RR ≈ 2.5). This is because we simulated a perfect,
infinitely large randomized trial. In this ideal scenario, where the
treated and untreated groups are exchangeable by design, the association
we measure is a consistent estimate of the causal effect.

\begin{center}\rule{0.5\linewidth}{0.5pt}\end{center}

\subsubsection{Solutions}\label{solutions}

\textless details\textgreater{}

\textless summary\textgreater Click to view
solutions\textless/summary\textgreater{}

\paragraph{Part A: Conceptual
Questions}\label{part-a-conceptual-questions-1}

\begin{enumerate}
\def\labelenumi{\arabic{enumi}.}
\item
  \textbf{Association vs.~Causation:} Let's consider the observed
  association between drinking coffee (\(A\)) and heart disease (\(Y\)).
  Many studies show coffee drinkers have a higher risk. However, coffee
  drinkers are also more likely to be smokers (\(L\)). Smoking is a
  known cause of heart disease. Therefore, an association between coffee
  and heart disease may appear even if coffee has no effect, simply
  because smoking is a \textbf{common cause} of both. The association is
  real, but it is \textbf{confounded} by smoking.

  \begin{itemize}
  \tightlist
  \item
    \textbf{Causal Diagram:}
  \end{itemize}
\item
  \textbf{Choosing a Measure:}

  \begin{itemize}
  \tightlist
  \item
    \textbf{Intervention 1:}

    \begin{itemize}
    \tightlist
    \item
      Risk in unexposed = 4/100 = 0.04. Risk in exposed = 2/100 = 0.02.
    \item
      \textbf{RR:} 0.02 / 0.04 = 0.5 (The risk is halved).
    \item
      \textbf{RD:} 0.02 - 0.04 = -0.02 (Prevents 2 cases per 100
      people).
    \end{itemize}
  \item
    \textbf{Intervention 2:}

    \begin{itemize}
    \tightlist
    \item
      Risk in unexposed = 40/100 = 0.40. Risk in exposed = 30/100 =
      0.30.
    \item
      \textbf{RR:} 0.30 / 0.40 = 0.75 (A 25\% risk reduction).
    \item
      \textbf{RD:} 0.30 - 0.40 = -0.10 (Prevents 10 cases per 100
      people).
    \end{itemize}
  \item
    \textbf{Interpretation:} The \textbf{Risk Difference (RD)} better
    reflects public health impact, as Intervention 2 prevents 5 times as
    many cases (10 vs.~2 per 100). The \textbf{Risk Ratio (RR)} better
    reflects the relative strength of the effect, as Intervention 1 cuts
    the risk in half, a stronger relative effect than the 25\% reduction
    from Intervention 2.
  \end{itemize}
\item
  \textbf{Counterfactuals:} It's impossible to know because we cannot
  observe the \textbf{counterfactual} outcome: what would have happened
  to \emph{that specific patient} at \emph{that specific time} if they
  had \emph{not} taken the aspirin. Maybe the headache was going to
  resolve on its own at that exact moment. This is the
  \textbf{fundamental problem of causal inference}.
\end{enumerate}

\paragraph{Part B: Computation from a
Table}\label{part-b-computation-from-a-table-1}

\begin{enumerate}
\def\labelenumi{\arabic{enumi}.}
\item
  \textbf{Risks:}

  \begin{itemize}
  \tightlist
  \item
    Risk in treated: \(Pr(Y=1|A=1) = 80 / 1000 = 0.08\) (or 8\%).
  \item
    Risk in untreated: \(Pr(Y=1|A=0) = 50 / 1000 = 0.05\) (or 5\%).
  \end{itemize}
\item
  \textbf{Risk Difference (RD):}

  \begin{itemize}
  \tightlist
  \item
    \(RD = 0.08 - 0.05 = 0.03\).
  \item
    \emph{Interpretation:} There are 3 additional cases of stroke for
    every 100 people treated with the drug compared to those not
    treated.
  \end{itemize}
\item
  \textbf{Risk Ratio (RR):}

  \begin{itemize}
  \tightlist
  \item
    \(RR = 0.08 / 0.05 = 1.6\).
  \item
    \emph{Interpretation:} People who take the drug are 1.6 times as
    likely to have a stroke as people who do not.
  \end{itemize}
\item
  \textbf{Odds Ratio (OR):}

  \begin{itemize}
  \tightlist
  \item
    Odds in treated: \((80/1000) / (920/1000) = 80/920 \approx 0.087\).
  \item
    Odds in untreated:
    \((50/1000) / (950/1000) = 50/950 \approx 0.053\).
  \item
    \(OR = (80/920) / (50/950) = (80 \times 950) / (920 \times 50) \approx 1.65\).
  \item
    \emph{Interpretation:} The odds of having a stroke are 65\% higher
    in the treated group compared to the untreated group. (Note that
    because the outcome is relatively rare, the OR of 1.65 is close to
    the RR of 1.6).
  \end{itemize}
\end{enumerate}

\textless/details\textgreater{}

\subsection{Appendix: Comprehension and
Application}\label{appendix-comprehension-and-application-1}

\subsubsection{Part A: Conceptual Questions
🧠}\label{part-a-conceptual-questions-2}

\begin{enumerate}
\def\labelenumi{\arabic{enumi}.}
\item
  \textbf{The Adage:}~Explain in your own words, using a simple public
  health example (other than lighters/smoking or coffee/heart disease),
  why ``association is not causation.'' Draw a simple causal diagram for
  your example.
\item
  \textbf{Choosing a Measure:}~An intervention reduces the risk of a
  disease from 4 in 100 to 2 in 100. Another intervention reduces the
  risk from 40 in 100 to 30 in 100.

  \begin{itemize}
  \item
    Calculate the Risk Ratio (RR) and Risk Difference (RD) for both
    interventions.
  \item
    Which measure (RR or RD) better reflects the public health impact
    (i.e., the number of cases prevented)? Which measure better reflects
    the strength of the effect relative to the baseline risk?
  \end{itemize}
\item
  \textbf{Counterfactuals:}~A patient took an aspirin and their headache
  went away. Why is it impossible to know for sure if the aspirin had a
  causal effect on that specific patient's headache? What is the
  specific term for this issue?
\end{enumerate}

\begin{center}\rule{0.5\linewidth}{0.5pt}\end{center}

\subsubsection{Part B: Computation from a Table
🔢}\label{part-b-computation-from-a-table-2}

Consider the following data from a hypothetical study of 2000
individuals on the effect of a new blood pressure drug (A) on the
one-year risk of stroke (Y).

\begin{longtable}[]{@{}llll@{}}
\toprule\noalign{}
& Stroke (Y=1) & No Stroke (Y=0) & Total \\
\midrule\noalign{}
\endhead
\bottomrule\noalign{}
\endlastfoot
\textbf{Drug (A=1)} & 80 & 920 & 1000 \\
\textbf{No Drug (A=0)} & 50 & 950 & 1000 \\
\end{longtable}

Calculate and interpret the following~\textbf{associational}~measures:

\begin{enumerate}
\def\labelenumi{\arabic{enumi}.}
\tightlist
\item
  Risk of stroke in the treated and untreated groups.
\item
  Risk Difference (RD).
\item
  Risk Ratio (RR).
\item
  Odds Ratio (OR).
\end{enumerate}

\begin{center}\rule{0.5\linewidth}{0.5pt}\end{center}

\subsubsection{Part C: Application with R
💻}\label{part-c-application-with-r-1}

Let's simulate a small dataset to see how this works in R.

\textbf{Scenario:}~We will simulate a population where a
treatment~\texttt{A}~has a true causal effect on a binary
outcome~\texttt{Y}. We'll start with a simple case: a randomized trial
where there is~\textbf{no confounding}.

\begin{verbatim}
# Load the tidyverse package for data manipulation and plotting
# If you don't have it, run: install.packages("tidyverse")
library(tidyverse)

# --- 1. Simulate a simple Randomized Controlled Trial (RCT) ---
set.seed(579) # for reproducibility

n <- 2000 # number of individuals in our study
# In an RCT, treatment A is assigned randomly (e.g., with 50% probability).
# It does not depend on any other characteristic of the individual.
treatment <- rbinom(n, 1, 0.5)

# The baseline risk of the outcome Y is 10%
# The true causal effect of the treatment is to increase the probability by 15 percentage points
# Pr(Y=1 | A=0) = 0.10
# Pr(Y=1 | A=1) = 0.10 + 0.15 = 0.25
# The true Causal Risk Ratio is 0.25 / 0.10 = 2.5
# The true Causal Risk Difference is 0.25 - 0.10 = 0.15
prob_y <- ifelse(treatment == 1, 0.25, 0.10)
outcome <- rbinom(n, 1, prob_y)

# Create a data frame (our dataset)
rct_data <- data.frame(A = treatment, Y = outcome)

# Let's inspect the first few rows
head(rct_data)

# --- 2. Calculate Associational Measures from the Data ---

# We can create a 2x2 table to summarize the data
rct_table <- table(rct_data$A, rct_data$Y)
print("2x2 Table (Rows: Treatment A, Columns: Outcome Y)")
print(rct_table)

# Let's calculate the measures
summary_stats <- rct_data %>%
  group_by(A) %>%
  summarise(
    n = n(),
    events = sum(Y),
    risk = mean(Y) # For a binary 0/1 outcome, mean() is the same as risk
  ) %>%
  ungroup()

print("Summary Statistics by Treatment Group:")
print(summary_stats)

# Extract risks for each group
risk_treated <- summary_stats$risk[summary_stats$A == 1]
risk_untreated <- summary_stats$risk[summary_stats$A == 0]

# Calculate Associational RD and RR
assoc_rd <- risk_treated - risk_untreated
assoc_rr <- risk_treated / risk_untreated

# Calculate Odds Ratio
odds_treated <- risk_treated / (1 - risk_treated)
odds_untreated <- risk_untreated / (1 - risk_untreated)
assoc_or <- odds_treated / odds_untreated

# --- 3. Interpretation ---

cat("\n--- Results from our Simulated RCT ---\n")
cat("Associational Risk Difference:", round(assoc_rd, 3), "\n")
cat("Associational Risk Ratio:", round(assoc_rr, 3), "\n")
cat("Associational Odds Ratio:", round(assoc_or, 3), "\n")
\end{verbatim}

\textbf{R Syntax Explanation:}

\begin{itemize}
\item
  \texttt{set.seed(579)}: This function ensures that every time you run
  the code, the ``random'' numbers generated will be the same. This is
  crucial for making our results reproducible.
\item
  \texttt{rbinom(n,\ 1,\ 0.5)}: This generates~\texttt{n}~random numbers
  from a binomial distribution. Here, it's like flipping a
  coin~\texttt{n}~times, where the probability of getting a 1 (heads) is
  0.5. This simulates our random treatment assignment.
\item
  \texttt{ifelse(test,\ yes,\ no)}: This is a conditional statement. It
  checks if~\texttt{treatment\ ==\ 1}. If true, it assigns the
  value~\texttt{0.25}; if false, it assigns~\texttt{0.10}. This is how
  we build the true causal effect into our simulation.
\item
  \texttt{\%\textgreater{}\%}: This is the ``pipe'' operator from
  the~\texttt{dplyr}~package (part of~\texttt{tidyverse}). It takes the
  output of the expression on its left and ``pipes'' it in as the first
  argument to the function on its right. It makes the code easier to
  read.
\item
  \texttt{group\_by(A)}: This tells R to perform the next set of
  operations separately for each group defined by the
  variable~\texttt{A}~(treated and untreated).
\item
  \texttt{summarise(…)}: This creates a new data frame with summary
  statistics.~\texttt{n()}~counts the rows,~\texttt{sum(Y)}counts the
  events, and~\texttt{mean(Y)}~calculates the risk.
\item
  \texttt{cat(…)}: This function is like~\texttt{print()}, but gives us
  more control over formatting the output text.
  The~\texttt{\textbackslash{}n}~creates a new line.
\end{itemize}

\textbf{Interpretation of R Output:}

When you run this code, you will see that the associational measures you
calculated from the simulated data are very close to the~\emph{true
causal effects}~we built into the simulation (Associational RD ≈ 0.15,
Associational RR ≈ 2.5). This is because we simulated a perfect,
infinitely large randomized trial. In this ideal scenario, where the
treated and untreated groups are exchangeable by design, the association
we measure is a consistent estimate of the causal effect.

\begin{center}\rule{0.5\linewidth}{0.5pt}\end{center}

\subsubsection{Solutions}\label{solutions-1}

Click to view solutions

\paragraph{\texorpdfstring{\textbf{Part A: Conceptual
Questions}}{Part A: Conceptual Questions}}\label{part-a-conceptual-questions-3}

\begin{enumerate}
\def\labelenumi{\arabic{enumi}.}
\item
  \textbf{Association vs.~Causation:}~Let's consider the observed
  association between owning a yellow car (A) and being in a car
  accident (Y). Data might show that owners of yellow cars have a lower
  risk of accidents. However, it's unlikely the yellow paint itself
  provides a physical shield. A more plausible explanation is that
  people who choose to buy a yellow car (A) may be more cautious drivers
  in general (L). This personality trait of being cautious (L) is
  a~\textbf{common cause}~of both choosing a yellow car and having a
  lower risk of an accident. The association is real, but it
  is~\textbf{confounded}~by driver caution.

  \begin{itemize}
  \tightlist
  \item
    \textbf{Causal Diagram:}~In this diagram, an unmeasured variable L
    (Cautious Personality) has an arrow pointing to A (Buys Yellow Car)
    and an arrow pointing to Y (Car Accident). There is no direct arrow
    from A to Y.
  \end{itemize}
\item
  \textbf{Choosing a Measure:}

  \begin{itemize}
  \item
    \textbf{Intervention 1:}

    \begin{itemize}
    \item
      Risk in unexposed = 4/100 = 0.04. Risk in exposed = 2/100 = 0.02.
    \item
      \textbf{RR:}~0.02 / 0.04 =~\textbf{0.5}~(The risk is halved).
    \item
      \textbf{RD:}~0.02 - 0.04 =~\textbf{-0.02}~(Prevents 2 cases per
      100 people).
    \end{itemize}
  \item
    \textbf{Intervention 2:}

    \begin{itemize}
    \item
      Risk in unexposed = 40/100 = 0.40. Risk in exposed = 30/100 =
      0.30.
    \item
      \textbf{RR:}~0.30 / 0.40 =~\textbf{0.75}~(A 25\% risk reduction).
    \item
      \textbf{RD:}~0.30 - 0.40 =~\textbf{-0.10}~(Prevents 10 cases per
      100 people).
    \end{itemize}
  \item
    \textbf{Interpretation:}~The~\textbf{Risk Difference (RD)}~better
    reflects public health impact, as Intervention 2 prevents 5 times as
    many cases (10 vs.~2 per 100). The~\textbf{Risk Ratio (RR)}~better
    reflects the relative strength of the effect, as Intervention 1 cuts
    the risk in half, a stronger relative effect than the 25\% reduction
    from Intervention 2.
  \end{itemize}
\item
  \textbf{Counterfactuals:}~It's impossible to know because we cannot
  observe the~\textbf{counterfactual}~outcome: what would have happened
  to~\emph{that specific patient}~at~\emph{that specific time}~if they
  had~\emph{not}~taken the aspirin. Maybe the headache was going to
  resolve on its own at that exact moment. This is
  the~\textbf{fundamental problem of causal inference}.
\end{enumerate}

\paragraph{\texorpdfstring{\textbf{Part B: Computation from a
Table}}{Part B: Computation from a Table}}\label{part-b-computation-from-a-table-3}

\begin{enumerate}
\def\labelenumi{\arabic{enumi}.}
\item
  \textbf{Risks:}

  \begin{itemize}
  \item
    Risk in treated:~Pr(Y=1∣A=1)=80/1000=0.08~(or 8\%).
  \item
    Risk in untreated:~Pr(Y=1∣A=0)=50/1000=0.05~(or 5\%).
  \end{itemize}
\item
  \textbf{Risk Difference (RD):}

  \begin{itemize}
  \item
    RD=0.08−0.05=0.03.
  \item
    \emph{Interpretation:}~There are 3 additional cases of stroke for
    every 100 people treated with the drug compared to those not
    treated.
  \end{itemize}
\item
  \textbf{Risk Ratio (RR):}

  \begin{itemize}
  \item
    RR=0.08/0.05=1.6.
  \item
    \emph{Interpretation:}~People who take the drug are 1.6 times as
    likely to have a stroke as people who do not.
  \end{itemize}
\item
  \textbf{Odds Ratio (OR):}

  \begin{itemize}
  \item
    Odds in treated:~(80/1000)/(920/1000)=80/920approx0.087.
  \item
    Odds in untreated:~(50/1000)/(950/1000)=50/950approx0.053.
  \item
    OR=(80/920)/(50/950)=(80times950)/(920times50)approx1.65.
  \item
    \emph{Interpretation:}~The odds of having a stroke are 65\% higher
    in the treated group compared to the untreated group. (Note that
    because the outcome is relatively rare, the OR of 1.65 is close to
    the RR of 1.6).
  \end{itemize}
\end{enumerate}

\paragraph{\texorpdfstring{\textbf{Part C: R Code Output and
Interpretation}}{Part C: R Code Output and Interpretation}}\label{part-c-r-code-output-and-interpretation}

Here is the output from running the R code:

\begin{verbatim}
> head(rct_data)
  A Y
1 1 0
2 1 1
3 0 0
4 1 0
5 1 0
6 0 0

> print("2x2 Table (Rows: Treatment A, Columns: Outcome Y)")
[1] "2x2 Table (Rows: Treatment A, Columns: Outcome Y)"
   
      0   1
  0 897 101
  1 753 249

> print("Summary Statistics by Treatment Group:")
[1] "Summary Statistics by Treatment Group:"
# A tibble: 2 × 4
      A     n events  risk
  <int> <int>  <dbl> <dbl>
1     0   998    101 0.101
2     1  1002    249 0.248

--- Results from our Simulated RCT ---
Associational Risk Difference: 0.147 
Associational Risk Ratio: 2.454 
Associational Odds Ratio: 2.89 
\end{verbatim}

\begin{itemize}
\tightlist
\item
  \textbf{Interpretation:}~The calculated associational risk difference
  is 0.147, which is very close to the true causal RD of 0.15. The
  associational risk ratio is 2.454, which is very close to the true
  causal RR of 2.5. This demonstrates that in an ideal randomized
  setting without confounding, the association we can calculate directly
  from the data is a good estimate of the causal effect.
\end{itemize}




\end{document}
